\documentclass[12pt]{letter}

\usepackage{amsmath,amssymb}
\usepackage{nopageno}
\usepackage{enumitem}
\usepackage[a4paper, total={7.5in, 11.25in}]{geometry}
\usepackage{tikz}
\usepackage{pgfplots}
\pgfplotsset{compat=1.18}

\usetikzlibrary{calc,3d,decorations.pathmorphing,decorations.pathreplacing,decorations.markings,snakes,arrows.meta,patterns}
\tikzset{snake it/.style={decorate, decoration=snake}}

\begin{document}

Name: \hrulefill \quad Neptun code: \hrulefill

\begin{center}
(KTXFI1EBNF) Physics I. Test~2 \\
May~18,~2026 
\end{center}

\bigskip

\begin{enumerate}
\setcounter{enumi}{0}

\item Estimate the average depth of the Indian Ocean if a tsunami takes 24 hours to travel from Jakarta to Mogadishu. (Note: The propagation speed of shallow-water waves is $v=\sqrt{gh}$. The distance is $17000$\,km. $g=9.81\,\mathrm{m/s^{2}}.$)

\hfill (10~points)

% ------------------------------------------------------------------------------------------------------------------------------------

\bigskip
\item A projectile is launched at an angle of $45^\circ$ with an initial velocity of $50\,\text{m/s}$ ($g=9.81\,\mathrm{m/s^{2}}$). 
Calculate:
\begin{enumerate}
    \item[(a)] The time to reach the peak of the trajectory.
    \item[(b)] The maximum height.
    \item[(c)] The horizontal range (impact distance).
    \item[(d)] The total time of flight.
\end{enumerate}     

\hfill (10~points)

% ------------------------------------------------------------------------------------------------------------------------------------

\bigskip
\item We place a copper piece with a mass of 100 gramms and a temperature of $40{}^{\circ}C$ into 500 gramms of water at $80{}^{\circ}C$. What will be the common temperature? The specific heat capacity of copper is $3.85 \cdot 10^{2}\frac{J}{kg\,K}$ and the specific heat capacity of water is $4.19 \cdot 10^{3}\frac{J}{kg\,K}$.

\hfill (10~points)

% ------------------------------------------------------------------------------------------------------------------------------------

\bigskip
\item  What is the entropy change if the volume of $ m = 14\,\mathrm{g} $ of nitrogen changes from $ V_1 = 30\,\mathrm{l} \text{\ to\ } V_2 = 120\,\mathrm{l},$ and its temperature changes from $T_1 = 27^\circ\mathrm{C} \text{\ to\ } T_2 = 327^\circ\mathrm{C}?$ The molar mass is $ M = 28\,\mathrm{g/mol}$, $R=8.31\,\mathrm{\frac{J}{mol\cdot K}}$. 

\medskip \hfill (10~points)

% ------------------------------------------------------------------------------------------------------------------------------------

\bigskip
\item Determine the acceleration of a proton ($m = 1.67 \cdot 10^{-27}\,\mathrm{kg}$, $e=1.6\cdot 10^{-19}\,\mathrm{C}$) in an electric field of strength $ E = 10000\,\mathrm{N/C}. $ How many times greater is this acceleration than the acceleration due to gravity?

\medskip \hfill (10~points)

\end{enumerate}

% ------------------------------------------------------------------------------------------------------------------------------------

Requirement for obtaining the course signature: achieving at least 40\% (20\,points) on the midterm exam.

Grades: 0--25 points: Fail (1), 26--30 points: Pass (2), 31--37 points: Satisfactory (3), 38--45 points:~Good~(4), 46--50 points: Excellent (5).

\rightline{Dr.~Gabor FACSKO}
\rightline{\textit{facsko.gabor@uni-obuda.hu}}

\vfill

\end{document}


